\textbf{Phần 1. Trắc nghiệm lựa chọn}
\setcounterpageref{ex}{0}

\begin{ex} %[3D2V10]
Tính giá trị của biểu thức: $B = 7 \times 8 - 42 : 7 + 16$.
\choice
{$64$}
{\True $66$}
{$68$}
{$62$}
\loigiai{
Thực hiện nhân chia trước: $B = 56 - 6 + 16 = 50 + 16 = 66$.
\textbf{Đáp án của câu hỏi là B}
}
\end{ex}

\begin{ex} %[3D2V10]
Một tuần lễ có $7$ ngày. Hỏi trong $5$ tuần lễ và $3$ ngày có tất cả bao nhiêu ngày?
\choice
{$35$ ngày}
{$36$ ngày}
{$40$ ngày}
{\True $38$ ngày}
\loigiai{
Số ngày trong 5 tuần là: $5 \times 7 = 35$ (ngày).\
Tổng số ngày là: $35 + 3 = 38$ (ngày).
\textbf{Đáp án của câu hỏi là D}
}
\end{ex}

\begin{ex} %[3D2V10]
Tìm $x$ biết: $x : 7 + 18 = 27$.
\choice
{\True $63$}
{$56$}
{$70$}
{$49$}
\loigiai{
Ta có: $x : 7 = 27 - 18 = 9 \Rightarrow x = 9 \times 7 = 63$.
\textbf{Đáp án của câu hỏi là A}
}
\end{ex}

\begin{ex} %[3D2V10]
Một cuộn dây thép dài $56\text{ m}$. Người ta cắt đi $\frac{1}{7}$ cuộn dây thép đó. Hỏi cuộn dây thép còn lại bao nhiêu mét?
\choice
{$8\text{ m}$}
{$50\text{ m}$}
{\True $48\text{ m}$}
{$42\text{ m}$}
\loigiai{
Đoạn dây thép đã cắt là: $56 : 7 = 8\text{ (m)}$.\
Đoạn dây còn lại dài: $56 - 8 = 48\text{ (m)}$.
\textbf{Đáp án của câu hỏi là C}
}
\end{ex}

\begin{ex} %[3D2V10]
Trong một phép chia cho $7$, thương bằng $9$ và số dư là số dư lớn nhất có thể có. Số bị chia của phép chia đó là
\choice
{$63$}
{\True $69$}
{$68$}
{$70$}
\loigiai{
Vì số chia là 7 nên số dư lớn nhất là 6.\
Số bị chia là: $7 \times 9 + 6 = 63 + 6 = 69$.
\textbf{Đáp án của câu hỏi là B}
}
\end{ex}

\begin{ex} %[3D2V10]
Tháng Chạp (tháng 12) có ngày mùng 1 là thứ Sáu. Hỏi ngày 22 tháng 12 cùng năm đó là thứ mấy?
\choice
{Thứ Bảy}
{Thứ Năm}
{Chủ Nhật}
{\True Thứ Sáu}
\loigiai{
Khoảng cách giữa ngày 22 và ngày 1 là: $22 - 1 = 21$ (ngày).\
Vì $21 : 7 = 3$ (tuần chẵn) nên ngày 22 cũng là thứ Sáu.
\textbf{Đáp án của câu hỏi là D}
}
\end{ex}

\begin{ex} %[3D2V10]
Tính nhanh: $7 \times 3 + 7 \times 4 + 7 \times 3$.
\choice
{\True $70$}
{$63$}
{$77$}
{$84$}
\loigiai{
Đặt thừa số chung: $7 \times (3 + 4 + 3) = 7 \times 10 = 70$.
\textbf{Đáp án của câu hỏi là A}
}
\end{ex}

\begin{ex} %[3D2V10]
Một mảnh đất hình chữ nhật có chiều rộng là $7\text{ m}$, chiều dài gấp $4$ lần chiều rộng. Chu vi của mảnh đất đó là
\choice
{$65\text{ m}$}
{$60\text{ m}$}
{\True $70\text{ m}$}
{$75\text{ m}$}
\loigiai{
Chiều dài mảnh đất là: $7 \times 4 = 28\text{ (m)}$.\
Chu vi là: $(28 + 7) \times 2 = 35 \times 2 = 70\text{ (m)}$.
\textbf{Đáp án của câu hỏi là C}
}
\end{ex}

\begin{ex} %[3D2V10]
Tìm $x$ biết: $(x - 15) \times 7 = 35$.
\choice
{$18$}
{\True $20$}
{$22$}
{$25$}
\loigiai{
Ta có: $x - 15 = 35 : 7 = 5 \Rightarrow x = 5 + 15 = 20$.
\textbf{Đáp án của câu hỏi là B}
}
\end{ex}

\begin{ex} %[3D2V10]
Mẹ mua $7$ vỉ trứng gà, mỗi vỉ có $6$ quả và mua thêm $15$ quả trứng vịt. Hỏi mẹ đã mua tất cả bao nhiêu quả trứng?
\choice
{$52$ quả}
{$55$ quả}
{$60$ quả}
{\True $57$ quả}
\loigiai{
Số trứng gà là: $7 \times 6 = 42$ quả.\
Tổng số trứng là: $42 + 15 = 57$ quả.
\textbf{Đáp án của câu hỏi là D}
}
\end{ex}

\begin{ex} %[3D2V10]
Cho dãy số cách đều: $7; 14; 21; 28; \dots$ Số hạng thứ 8 của dãy số là
\choice
{\True $56$}
{$49$}
{$63$}
{$54$}
\loigiai{
Số hạng thứ 8 là: $7 \times 8 = 56$.
\textbf{Đáp án của câu hỏi là A}
}
\end{ex}

\begin{ex} %[3D2V10]
Bác Nam có $63$ cây cam trồng đều thành 7 hàng. Bác bán đi 2 hàng cam. Hỏi bác Nam còn lại bao nhiêu cây cam?
\choice
{$40$ cây cam}
{$42$ cây cam}
{\True $45$ cây cam}
{$48$ cây cam}
\loigiai{
Mỗi hàng có số cây cam là: $63 : 7 = 9$ (cây).\
Số hàng còn lại là: $7 - 2 = 5$ (hàng).\
Số cây cam còn lại là: $5 \times 9 = 45$ (cây).
\textbf{Đáp án của câu hỏi là C}
}
\end{ex}

\textbf{Phần 2. Trắc nghiệm đúng sai}
\setcounterpageref{ex}{0}

\begin{ex} %[3D2V10]
Xét tính đúng sai của các nhận định sau về bảng nhân và chia 7:
\choiceTF
{\True Một tuần lễ có đúng 7 ngày.}
{Trong phép chia cho 7, số dư lớn nhất có thể có là 7.}
{\True Nếu lấy tích của 7 với một số rồi chia cho 7 thì được chính số đó.}
{\True Tổng $7 \times 4 + 7 \times 6 = 70$.}
\loigiai{
\begin{itemize}
    \item Ý a) Đúng: Quy ước thời gian 1 tuần = 7 ngày.
    \item Ý b) Sai: Số dư lớn nhất là 6.
    \item Ý c) Đúng: $7 \times a : 7 = a$.
    \item Ý d) Đúng: $7 \times (4 + 6) = 7 \times 10 = 70$.
\end{itemize}
\begin{itemize}
    \item a) \textbf{Đúng}.
    \item b) \textbf{Sai}.
    \item c) \textbf{Đúng}.
    \item d) \textbf{Đúng}.
\end{itemize}
}
\end{ex}

\begin{ex} %[3D2V10]
Một lớp học có 35 học sinh chia đều thành 7 nhóm học tập:
\choiceTF
{\True Mỗi nhóm học tập có 5 bạn học sinh.}
{\True Nếu 3 nhóm cử đi trực nhật thì có 15 bạn trực nhật.}
{Số học sinh của 4 nhóm là 25 bạn.}
{\True Tổng số học sinh gấp 7 lần số học sinh của mỗi nhóm.}
\loigiai{
\begin{itemize}
    \item Ý a) Đúng: $35 : 7 = 5$ bạn.
    \item Ý b) Đúng: $3 \times 5 = 15$ bạn.
    \item Ý c) Sai: $4 \times 5 = 20$ bạn.
    \item Ý d) Đúng: $35 : 5 = 7$ lần.
\end{itemize}
\begin{itemize}
    \item a) \textbf{Đúng}.
    \item b) \textbf{Đúng}.
    \item c) \textbf{Sai}.
    \item d) \textbf{Đúng}.
\end{itemize}
}
\end{ex}

\textbf{Phần 3. Trả lời ngắn}
\setcounterpageref{ex}{0}

\begin{ex} %[3D2V10]
Tính giá trị biểu thức: $7 \times 9 - 35 : 7$.
\shortans{58}
\loigiai{
\textbf{Đáp số:} $58$\\
$7 \times 9 - 35 : 7 = 63 - 5 = 58$.
}
\end{ex}

\begin{ex} %[3D2V10]
Tìm số bị chia trong phép chia cho 7, biết thương là 8 và số dư là 5.
\shortans{61}
\loigiai{
\textbf{Đáp số:} $61$\\
$7 \times 8 + 5 = 56 + 5 = 61$.
}
\end{ex}

\begin{ex} %[3D2V10]
Hôm nay là thứ Ba ngày 5. Hỏi thứ Ba tuần sau là ngày mấy?
\shortans{12}
\loigiai{
\textbf{Đáp số:} $12$\\
Thứ Ba tuần sau là ngày: $5 + 7 = 12$.
}
\end{ex}

\begin{ex} %[3D2V10]
Một bao xi măng nặng 49 kg được chia đều vào 7 túi nhỏ. Hỏi 4 túi nhỏ như thế nặng bao nhiêu ki-lô-gam?
\shortans{28}
\loigiai{
\textbf{Đáp số:} $28$\\
Mỗi túi nặng: $49 : 7 = 7\text{ kg}$. 4 túi nặng: $4 \times 7 = 28\text{ kg}$.
}
\end{ex}

\begin{ex} %[3D2V10]
Tìm $x$ biết: $7 \times x = 90 - 41$.
\shortans{7}
\loigiai{
\textbf{Đáp số:} $7$\\
$7 \times x = 49 \Rightarrow x = 49 : 7 = 7$.
}
\end{ex}

\begin{ex} %[3D2V10]
Tính nhanh: $7 \times 2 + 7 \times 3 + 7 \times 5$.
\shortans{70}
\loigiai{
\textbf{Đáp số:} $70$\\
$7 \times (2 + 3 + 5) = 7 \times 10 = 70$.
}
\end{ex}
